\chapter{Construction specification and measured components}\label{ch:construction}
\section{An operation must preserve the study's meaning}
The implementation should make the scientific argument reproducible across
execution environments. A scientist who requests an average effect on the
original eligible cohort should receive that target whether the job runs in
a notebook, a service worker or a distributed warehouse. Removing closed
accounts, changing fold boundaries or averaging branch coefficients with
different weights changes the analysis even if every local function succeeds.
The unit of composition is therefore a versioned study and its artifacts.

The proposal supplies two small working components: historical event
reconstruction and crossed-cluster score reduction. It also supplies fitted
research studies and four JSON contracts. These establish selected semantics
and calculations. The general study compiler, service endpoints, user
workspaces, distributed scheduler and sibling adapters described here remain
proposed construction. Keeping that boundary explicit lets an implementation
team identify its actual starting point.

\section{Package boundaries and owned invariants}
The proposed packages divide responsibility by the invariant they preserve.
The data package owns historical identity, observation status, availability
and reconciled economic components. The design package owns the estimand,
intervention versions, identification route, assumptions and target weights.
The execution package owns folds, model dependencies, resource limits and
retry semantics. Estimation adapters own the translation to an explicitly
supported external method. The inference package owns score meaning,
dependence, resampling and simultaneous contrasts. The evidence package owns
immutable reconstruction, source identity and permission-aware exports.

These packages should exchange documented arrays and records rather than
opaque framework objects. A trajectory artifact records trajectory key,
calendar key, decision index, action and eligibility mask, pre-action
information, outcomes, censoring and model-training role. An irregular
history cannot be padded into a balanced panel without recording which
positions are structural absence, unobserved visits or actual zeros.
An adapter requiring equal-length histories must reject or explicitly
transform an incompatible artifact.

\begin{longtable}{P{31mm}P{54mm}P{57mm}}
\toprule Operation & Required input & Returned artifact or refusal\\\midrule\endhead
Reconstruct snapshot & Decision episodes, event revisions, entity and coverage rules & Selected revisions, features, observation masks, reconciliation\\
Compile design & Study, snapshot, graph assumptions, target population & Estimand and identification status, fold plan, supported methods\\
Fit nuisance/model & Training partition, method/version, seed, resource limit & Model, training lineage, convergence and support diagnostics\\
Evaluate scores & Frozen models, evaluation partition, policy/target weights & Defined score vector, cluster keys, extreme-contribution diagnostics\\
Infer contrasts & Complete score partitions, covariance specification & Point estimate, labeled uncertainty, covariance and limitations\\
Export evidence & Completed run and authorized artifact references & Reconstruction manifest, scientific report and financial account\\
\bottomrule
\caption{Proposed operation contracts. Each refusal identifies a specific
missing condition; it is not a request to silently choose a new target.}
\label{tab:operations}
\end{longtable}

A model fit can finish while its transition parameters remain weakly
determined. A design can be valid while a particular fitted model reaches
its iteration cap. A job can be canceled before a complete covariance is
available. The result should therefore carry separate computational,
scientific and evidence states. Suggested computational states are queued,
running, complete, failed and canceled. Scientific states distinguish
identified under declared assumptions, partially identified, conditional
scenario and unresolved. Evidence states distinguish exact construction,
synthetic study, retrospective bank study and prospective evaluation.
One green success flag cannot communicate all three.

\section{The implemented temporal reference}
The small Python package is named \texttt{causalbank-research}, version 0.2.0.
Its \texttt{Event} and \texttt{Decision} records and
\texttt{compile\_snapshot} function implement the revision ordering in
Chapter~\ref{ch:data}. For example:
\begin{lstlisting}[language=Python]
from causalbank import Event, Decision, compile_snapshot

events = [
    Event("credit", "H1", "2026-01-02T00:00:00Z",
          "2026-01-03T00:00:00Z", 1, 100.0, "USD"),
    Event("credit", "H1", "2026-01-02T00:00:00Z",
          "2026-01-10T00:00:00Z", 2, 60.0, "USD"),
]
decisions = [Decision("d1", "H1", "2026-01-08T00:00:00Z")]
snapshot = compile_snapshot(
    events, decisions, lookback_days=30,
    currency="USD", feed_complete=True)
assert snapshot["rows"][0]["amount"] == 100.0
\end{lstlisting}
The result lists the selected event revision and a hash of the declared
input account. Reversing input event order preserves the result. A duplicate
event revision, changed event owner, nonfinite amount, incompatible currency
or timestamp without a timezone produces an error. The feed-completeness
declaration is required before absence becomes an additive zero.

This reference deliberately uses an inspectable per-decision algorithm.
It is not the implementation to scan a banking group's full transaction
history. A warehouse adapter should implement the same semantics with
indexed intervals, partition pruning and revision-aware joins, then pass the
same independent account fixture. Hashing a complete local input collection
also differs from production content addressing: a large job should use
immutable dataset manifests and partition digests instead of rereading every
byte at each operation.

Availability at exactly the recorded decision instant is treated as known
in the reference. A real decision system may require a strict predecessor
rule or a publication-latency margin. The source sequence number and the
decision system's clock precision determine that choice. It must be a
versioned convention; changing it creates a new snapshot.

\section{The implemented score reducer}
Let $x_j\in\mathbb R^d$ be an already computed, equally weighted score
vector. Write $n$ for its row count, $\bar x=n^{-1}\sum_jx_j$, and
$S_g=\sum_{j\in g}x_j$, $n_g=|g|$ for a declared cluster. The centered
cluster contribution is $U_g=S_g-n_g\bar x$. For $G$ nonempty clusters,
the component covariance used by the reference is
\[
 \widehat V_G=\frac{G}{G-1}\frac1{n^2}\sum_gU_gU_g'.
\]
Entity, date and entity--date intersections give
$\widehat V=\widehat V_E+\widehat V_T-\widehat V_{ET}$.
Retaining intersection groups matters when multiple scores occupy the same
entity/date cell. Replacing the last term by an ordinary row covariance
would then generally compute a different quantity.

Raw counts and sums are additive across partitions. The reducer merges
those sufficient statistics and centers once with the global mean.
Centering each partition and averaging its covariance generally loses
between-partition terms. The supplied check compares the merge result with
an independently centered batch calculation on a crossed 63-row fixture;
maximum discrepancies are below $2\times10^{-16}$ in this environment.

\begin{lstlisting}[language=Python]
from causalbank import ScoreReducer

left = ScoreReducer(n_entities=10000, n_dates=120, dimension=2)
right = ScoreReducer(n_entities=10000, n_dates=120, dimension=2)
# Each worker adds disjoint rows of precomputed scores and integer keys.
left.add(score_part_a, entity_part_a, date_part_a)
right.add(score_part_b, entity_part_b, date_part_b)
result = left.merge(right).result()
\end{lstlisting}
The caller owns score construction, target weighting and the sampling
assumptions. A raw regression coefficient is not automatically a score.
For a nonlinear estimating equation, its Jacobian and normalization must
be applied consistently. Unequal target weights need a defined weighted
score and denominator; they are not supplied by this prototype.

The covariance can be indefinite in a finite sample. The reference reports
its minimum eigenvalue and a flag instead of silently forcing a positive
matrix. A consumer must then use a scientifically justified alternative or
withhold the affected interval. The component also rejects nonfinite
scores, incompatible shapes and too few nonempty clusters. It does not
choose a bandwidth, bootstrap method or independent-unit definition for
the analyst.

\section{What the scale measurement actually times}
The scale fixture creates two-dimensional synthetic scores in chunks of
100,000 rows, with 10,000 declared entity keys and 120 date keys. Three
fresh child processes run at each size. The total timer includes score
generation, accumulation and final covariance; the reduction timer isolates
accumulation. Linux peak resident memory includes the interpreter and
loaded libraries.

\begin{table}[htbp]\centering\small
\begin{tabular}{@{}rrrr@{}}
\toprule Rows & Total seconds & Reduce seconds & Peak RSS (MiB)\\\midrule
100,000 & 0.097 & 0.078 & 77.1\\
1,000,000 & 0.233 & 0.133 & 115.2\\
10,000,000 & 1.186 & 0.630 & 123.1\\
\bottomrule\end{tabular}

\caption{Median of three local runs per size. Precomputed-score reduction
only: no storage reads, temporal reconstruction, nuisance fitting, causal
graph search or POMDP parameter fitting is included.}
\label{tab:scale-results}
\end{table}

The measured ten-million-row operation takes about 1.2 seconds and
123 MiB peak resident memory in this environment. These numbers show
that a bounded score reduction can avoid retaining all rows. They cannot
be advertised as the latency of a ten-million-row causal analysis.
The dense intersection arrays require
$O((E+T+ET)d)$ score storage plus counts. With ten million entities,
even 120 dates makes that representation inappropriate. A sparse grouped
or warehouse aggregation is then required and must be measured separately.
The current measurement does not establish a universal streaming-memory
bound independent of cluster cardinality.

The benchmark records Python and NumPy versions, platform and component
hashes. Shared-host load can change timing. There was no dedicated machine
reservation, storage benchmark or comparison with Ray. A production
qualification should specify cores, memory limits, execution concurrency,
data format, warm/cold caches and complete stage costs before committing
to a service objective.

\section{A reproducible job is a dependency graph}
A compiled job should reference a frozen snapshot and a deterministic fold
assignment. A nuisance fit depends only on its training role. Score tasks
depend on the exact model versions and their evaluation partitions.
Inference depends on the complete set of score partitions, with declared
target weights and covariance method. Report generation depends on the
financial account and both scientific and computational statuses.

The artifact identity should hash canonical configuration, input manifest
digests, code/dependency lock identity and random seeds. A model upgrade,
changed treatment mapping or later outcome snapshot changes that identity.
An authorization decision remains dynamic even when content is immutable:
a previously readable artifact may cease to be exportable after a permission
change. Content identity must not be confused with access authority.

Retries require an exactly-once contribution rule. A failed worker may
write the same score partition twice; the manifest must select one committed
artifact for that partition, rather than summing both. Committing the
partition manifest atomically is preferable to inferring completion from
the presence of a partially written file. A reducer starts only when all
expected partitions are present and valid.

Checkpointing is method-specific. An EM checkpoint includes model
parameters, iteration, objective and sufficient-statistic configuration.
A bootstrap checkpoint includes resample identity and fit diagnostics.
A graph-search checkpoint includes candidate constraints and conditioning
tests already performed. Resuming with changed data or constraints is a
new run. Saving an arbitrary Python process state is not the public
reconstruction contract.

\section{Three concrete repository adoption journeys}
FinMarketing can export assignment episodes, decision-time covariates and
a group contribution ledger. The proposed integration reconstructs the
snapshot, computes cohort effects and returns joint balance/contribution
uncertainty with an explicit target and sign. Acceptance compares the same
study with FinMarketing's existing calculation and reconciles transfers
account by account.

dataml can export a recovery ensemble with observed/missing masks, the
measurement model, draw identifiers and training partition provenance.
The causal consumer repeats the declared analysis under compatible draws
and reports recovery sensitivity. The initial integration should use an
exact small corrupted-data fixture whose observed entries remain unchanged.
It fails if a recovered outcome is exposed as a historical policy feature
or uncertainty is lost when the ensemble becomes a dataframe.

BayesFilter can supply a filtering or likelihood component through a narrow
numerical interface. A controlled model specifies whether the action factor
enters its historical posterior and when the reward is observed. Its
probabilities and gradients are compared with exact short-path sums before
larger fits. The causal library retains ownership of the intervention
interpretation. MacroFinance can similarly consume response-path covariance
and shock restrictions; cardnpv additionally requires an explicit contrast
direction conversion. These are proposed acceptance journeys, not existing
inter-repository calls.

\section{Building the independent application}
The standalone workbench should guide the bank user through subject and
action selection, eligibility, the decision-time account, design declaration,
method comparison and economic outcomes. Discovery should display ambiguous
endpoints and measurements or interventions that could distinguish them.
Latent-state analysis should show filtering uncertainty, fit diagnostics
and conditional policy value, preserving its identification restrictions.

The proposed SDK and evidence export should expose the operation records.
The interface should explain unsupported requests in the user's terms:
no untreated branches remain after the last rollout date, or severe-state
assistance was never observed. Worker identifiers and dependency hashes
belong in expandable technical evidence.

A bank deployment should reuse identity, catalog, warehouse and storage
services. Workers execute near authorized data; the control service holds
study metadata and permitted references. Regional boundaries cover derived
scores and fitted models as well as rows. R, Java and Python adapters each
require environment qualification and applicable license review.

Release acceptance reconstructs an exported study in a clean qualified
environment; detects late-data, support and sign errors; resumes interrupted
jobs without duplicate contributions; and reconciles financial totals
independently. Review-effort savings must then be measured on equivalent tasks.
